%%
%% Kapitel 3 -- Stand der Technik: datengetriebene Systemidentifikation
%%
%%%%%%%%%%%%%%%%%%%%%%%%%%%%%%%%%%%%%%%%%%%%%%%%%%%%%%%%%%%%%%%%%%%%%

\chapter{Stand der Technik: datengetriebene Systemidentifikation}
\label{ch.standderTechnik}

Die Modellierung dynamischer Systeme aus Messdaten lässt sich in drei
Klassen einteilen \cite{isermannAutomotiveControlModeling2022}:

\begin{itemize}
  \item \textbf{White-Box-Modelle} basieren vollständig auf den
    Erhaltungsgleichungen für Masse, Energie und Impuls. Vertreter für
    Verbrennungsmotoren sind Mittelwertmodelle (Mean-Value Engine
    Models, MVEM) und nulldimensionale Brennraummodelle
    \cite{guzzellaIntroductionModelingControl2010, chmelaNulldimensionaleModellierung2019}.
    Ihre Stärke ist die physikalische Interpretierbarkeit und
    Extrapolierbarkeit; ihre Schwäche der hohe Aufwand für die
    Parameteridentifikation sowie Modellungenauigkeiten in Komponenten
    wie Kühlern, die nur eingeschränkt physikalisch parametrisierbar
    sind \cite{tangAccuracyMeanValue2024}.
  \item \textbf{Black-Box-Modelle} verwenden ausschließlich
    Eingangs-Ausgangs-Daten ohne Anspruch auf physikalische
    Interpretierbarkeit. Klassische Vertreter sind ARX, ARMAX und OE,
    im nichtlinearen Fall neuronale Netze.
  \item \textbf{Grey-Box- bzw.\ Hybridmodelle} kombinieren beide
    Ansätze, indem sie physikalische Modellstrukturen mit
    datengetriebenen Parametern verbinden
    \cite{wulffPhysicsinformedRegressionModels2025, ellisMachineLearningEnhanced2021, subramanianWhiteboxMachineLearning2021}.
    Die in dieser Arbeit eingesetzten Verfahren \sindy{} und \edmd{}
    \cite{bruntonDiscoveringGoverningEquations2016, williamsDataDrivenApproximation2015}
    sind hier einzuordnen: Sie werden zwar datengetrieben aus Messdaten
    geschätzt, liefern aber (anders als ein neuronales Netz) eine
    explizite Gleichung mit einsehbaren Termen und Koeffizienten. Mit
    einer physikalisch motivierten Bibliothek
    (Abschnitt~\ref{ssec.sot.sindy.bib}) werden sie zu vollwertigen
    Grey-Box-Modellen.
\end{itemize}

Die vorliegende Arbeit bewegt sich im Black-Box- und Grey-Box-Bereich.
Formal gehören alle hier betrachteten datengetriebenen Verfahren zum
\emph{überwachten maschinellen Lernen}: Aus gemessenen Ein- und
Ausgangsgrößen wird eine Abbildung geschätzt, die den zeitlichen
Fortgang des Systems vorhersagt. Ihnen ist ein Kern gemeinsam: Sie
verschieben die gemessene Zeitreihe um einen Schritt (bzw.\ bilden ihre
Zeitableitung) und bestimmen anschließend diejenigen Koeffizienten, die
diese zeitliche Veränderung im Sinne kleinster Quadrate am besten
erklären. 


\section{Regression der kleinsten Quadrate}
\label{sec.sot.lsq}

Vor der getrennten Betrachtung der einzelnen Verfahren lohnt der Blick
auf ihren gemeinsamen Kern, die Regression. Sie wertet den
statistischen Zusammenhang zwischen gemessenen Größen aus. 
Alle in diesem Kapitel diskutierten
datengetriebenen Methoden bestimmen eine Koeffizientenmatrix
$\boldsymbol\Xi$, die eine Zielgröße $\mathbf Y$ möglichst gut aus einer
Matrix von Merkmalen (Features) $\boldsymbol\Theta$ rekonstruiert,

\begin{equation}
  \label{eq.lsq}
  \mathbf Y \;=\; \boldsymbol\Theta\,\boldsymbol\Xi + \mathbf E,
  \qquad
  \hat{\boldsymbol\Xi} \;=\; \argmin_{\boldsymbol\Xi}\,
  \bigl\lVert \mathbf Y - \boldsymbol\Theta\,\boldsymbol\Xi \bigr\rVert_2^2 .
\end{equation}


Bei vollem Spaltenrang von $\boldsymbol\Theta$ hat
(\ref{eq.lsq}) die geschlossene Lösung
$\hat{\boldsymbol\Xi} = \boldsymbol\Theta^{+}\mathbf Y$ mit der
Moore-Penrose-Pseudoinversen
$\boldsymbol\Theta^{+} =
(\boldsymbol\Theta^{\!\top}\boldsymbol\Theta)^{-1}\boldsymbol\Theta^{\!\top}$.
Geometrisch ist die kleinste-Quadrate-Schätzung
$\hat{\mathbf Y} = \boldsymbol\Theta\hat{\boldsymbol\Xi}$ die
Orthogonalprojektion von $\mathbf Y$ auf den von den Spalten von
$\boldsymbol\Theta$ aufgespannten Unterraum; das Residuum
$\mathbf E = \mathbf Y - \hat{\mathbf Y}$ steht senkrecht darauf
(Bild~\ref{fig.lsq}).

\begin{figure}[bht]
  \centering
  \includegraphics[width=.62\textwidth]{lsq_projektion}
  \caption{Kleinste-Quadrate-Schätzung als Orthogonalprojektion. Die
    Spalten der Merkmalsmatrix $\boldsymbol\Theta$ (angedeutet durch
    $\boldsymbol\theta_1,\boldsymbol\theta_2$) spannen einen Unterraum
    auf. Die Schätzung $\hat{\mathbf y}=\boldsymbol\Theta\hat{\boldsymbol\Xi}$
    ist die orthogonale Projektion der Messung $\mathbf y$ in diesen
    Unterraum, das Residuum $\mathbf e = \mathbf y-\hat{\mathbf y}$ steht
    senkrecht darauf. Alle datengetriebenen Verfahren dieses Kapitels
    instanziieren dieses Grundbild mit unterschiedlichem
    $\boldsymbol\Theta$.}
  \label{fig.lsq}
\end{figure}

Alle folgenden Verfahren basieren auf dieser Regression der kleinsten Quadrate. Sie unterscheiden sich allein darin, was als
Merkmal in die Regression eingeht, welche Größe vorhergesagt wird
und \emph{ob} die Lösung dünn besetzt erzwungen wird.

%%Kollinearität

Aus (\ref{eq.lsq}) folgen zugleich die Grenzen des Ansatzes. Sind Spalten
von $\boldsymbol\Theta$ linear voneinander abhängig oder nahezu abhängig,
ist $\boldsymbol\Theta^{\!\top}\boldsymbol\Theta$ schlecht konditioniert.
Eine Lösung liefert die Pseudoinverse weiterhin. Deren Koeffizienten
reagieren jedoch empfindlich auf kleine Änderungen der Daten, verteilen den
Beitrag zweier ähnlicher Spalten willkürlich und nehmen dabei auch falsche
Vorzeichen an. Das Residuum bleibt klein, die physikalische Lesart der
Gleichung geht verloren. Als Maß dieser Multikollinearität dient der
Varianzinflationsfaktor (VIF) der einzelnen Spalte. Breite Bibliotheken aus
Polynom- und Kreuztermen verschärfen den Effekt, weil die Spaltenzahl mit
dem Polynomgrad und der Zahl der Zustände kombinatorisch wächst. Zwei
Gegenmittel greifen unmittelbar in die Minimierung nach (\ref{eq.lsq}) ein:
Die Ridge-Regression bestraft die quadrierte Koeffizientennorm und
stabilisiert die Inversion, die Lasso-Regression bestraft die Betragsnorm
und drängt schwache Terme auf null. Beide Wege kehren in den Verfahren
dieses Kapitels als Regularisierung und als Schwellenwertauswahl wieder
(Abschnitt~\ref{ssec.sot.sindy.robust}).

%%Normierung

Zwei weitere Grenzen liegen in den Daten und nicht im Schätzer. Die Spalten
von $\boldsymbol\Theta$ tragen verschiedene Einheiten und Größenordnungen.
Ohne Zentrierung und Normierung bestimmen die betragsgroßen Spalten die
Lösung, und die Beiträge kleiner physikalischer Größen fallen aus der
Auswahl. Abhilfe schafft eine Spaltennormierung. Bei der Z-Normierung
wird jede Spalte um ihren Mittelwert zentriert und durch ihre empirische
Standardabweichung geteilt,
$\tilde{\boldsymbol\theta}_j = (\boldsymbol\theta_j - \mu_j)/\sigma_j$.
Für Regressionsbibliotheken ist die reine Skalierung ohne Zentrierung
üblich, weil die Zentrierung an der konstanten Spalte scheitert
($\sigma_j = 0$) und nach der Rückrechnung einen zusätzlichen
Konstantterm erzeugt. Geschätzt wird auf den normierten Spalten,
anschließend werden die Koeffizienten in die ursprünglichen Einheiten
zurückgerechnet. Für schwellwertbasierte Auswahlverfahren ist dieser
Schritt Voraussetzung, weil ein gemeinsamer Schwellwert vergleichbare
Größenordnungen aller Spalten verlangt.

Ähnlich wirkt der Abtasttakt. Ist er gegenüber der Prozessdynamik
zu fein, ändert sich der Zustand zwischen zwei Abtastwerten kaum, die
numerische Ableitung geht gegen null und die Regression passt sich dem
Sensorrauschen an. Ist er zu grob, fehlen die schnellen Anteile der Dynamik
und die geschätzte Ableitung ist verzerrt. Hinzu kommt die Überanpassung,
denn mit jeder weiteren Spalte sinkt das Residuum $\mathbf E$, ohne dass die
Prognose auf ungesehenen Daten besser wird. Regressionsmodelle interpolieren
innerhalb des im Training abgedeckten Betriebsbereichs; außerhalb davon
verlieren besonders Polynombibliotheken hohen Grades ihre Gültigkeit. Die
aus der Rückführung stammenden Grenzen sind davon zu trennen.


\section{Lineare Zeitreihenmodelle als Baseline: ARX und NARX}
\label{sec.sot.arx}

Die linearen Polynommodelle bilden den methodischen Anker jeder
datengetriebenen Modellierung \cite[Kap.~6]{isermannAutomotiveControlModeling2022}.
Das ARX-Modell

\begin{equation}
  \label{eq.arx}
  A(q)\,y(k) \;=\; B(q)\,u(k) + e(k)
\end{equation}

verknüpft den Ausgang $y(k)$ mit dem Eingang $u(k)$ und dem
Gleichungsfehler $e(k)$. Darin sind
$A(q) = 1 + a_1 q^{-1} + \dots + a_{n_a} q^{-n_a}$ und
$B(q) = q^{-n_k}\bigl(b_1 + b_2 q^{-1} + \dots + b_{n_b} q^{-(n_b-1)}\bigr)$
Polynome im Rückwärtsverschiebungsoperator $q^{-1}$ mit
$q^{-1}y(k) = y(k-1)$. Die Ordnungen $n_a$ und $n_b$ legen fest, wie
viele vergangene Aus- und Eingangswerte eingehen; die Totzeit $n_k$
gibt an, nach wie vielen Abtastschritten der Eingang zu wirken beginnt. Das Modell
ist linear in den Parametern und damit die unmittelbarste Instanz von
(\ref{eq.lsq}): Die Merkmalsmatrix $\boldsymbol\Theta$ enthält die
verzögerten Aus- und Eingangswerte, die Zielgröße ist $y(k)$. Die
Schätzung per kleinster Quadrate ist konsistent, sobald der
Gleichungsfehler $e(k)$ weiß ist, das Störmodell $1/A(q)$ also mit in
der Modellklasse liegt, und das Eingangssignal $u(k)$ persistente
Anregung ausreichender Ordnung besitzt
\cite[Kap.~8--9]{isermannAutomotiveControlModeling2022}. Im
geschlossenen Kreis ist die erste Bedingung die kritische. ARMAX erweitert (\ref{eq.arx}) um ein
gleitendes Mittel des Rauschens, OE trennt System- und Rauschmodell
vollständig.

Die Linearität ist eine harte Annahme. Die in
Abschnitt~\ref{ssec.grundlagen.luftpfad.kennzahlen} hergeleitete
Nichtlinearität in $\dot m_L$ und $\TLT$ kann ARX nur lokal abbilden.
Zudem erzeugt ARX keine interpretierbare Modellgleichung im Sinne
physikalischer Terme, sondern eine Übertragungsfunktion in
Verschiebungsoperatoren, deren Koeffizienten den Linearisierungspunkt
implizit kodieren. Daraus folgt eine schwache Übertragbarkeit zwischen
Betriebspunkten oder Motoren \cite{atamAdvancedAirPath2018}. Aus diesen
Gründen dient ARX in dieser Arbeit als Baseline, nicht als Zielmodell.

Als nichtlineare Erweiterung existiert die NARX-/NARMAX-Familie, bei der
die linearen Polynome durch nichtlineare Funktionen des verzögerten
Aus- und Eingangs ersetzt werden. \cite{aranDieselEngineAirpath2020}
verwendet eine MISO-NFIR-Variante für die Diesel-Luftpfad-Identifikation.
Methodisch lässt sich \sindy{} als sparsifizierte NARX-Variante
auffassen, in der eine vordefinierte nichtlineare Bibliothek anstelle
einer freien Funktionsklasse verwendet wird. Diese Lesart ordnet
\sindy{} in das klassische Identifikationsvokabular ein, ändert aber
nichts an den oben benannten Limitierungen der linearen Modellklasse.


%%%%%%%%%%%%%%%%%%%%%%%%%%%%%%%%%%%%%%%%%%%%%%%%%%%%%%%%%%%%%%%%%%%%%
%%%%%%%%%%%%%%%%%%%%%%%%%%%%%%%%%%%%%%%%%%%%%%%%%%%%%%%%%%%%%%%%%%%%%
%%  Dynamic Mode Decomposition
%%%%%%%%%%%%%%%%%%%%%%%%%%%%%%%%%%%%%%%%%%%%%%%%%%%%%%%%%%%%%%%%%%%%%
%%%%%%%%%%%%%%%%%%%%%%%%%%%%%%%%%%%%%%%%%%%%%%%%%%%%%%%%%%%%%%%%%%%%%


\section{Dynamic Mode Decomposition (DMD)}
\label{sec.sot.edmd}

DMD ist eng mit dem Koopman-Operator verbunden. Der Koopman-Operator wirkt linear auf einen unendlichdimsensionalen Raum

Der Koopman-Operator $\mathcal K$ wirkt linear auf den
unendlichdimensionalen Raum der Observablen einer dynamischen Abbildung
\cite{mezicSpectralPropertiesDynamical2005}. Für eine Observable
$\psi: \mathcal X \to \mathbb C$ und eine diskrete Dynamik
$\mathbf x_{k+1} = F(\mathbf x_k)$ gilt

\begin{equation}
  \label{eq.koopman}
  (\mathcal K \psi)(\mathbf x_k) \;=\; \psi(F(\mathbf x_k)) \;=\; \psi(\mathbf x_{k+1}).
\end{equation}

Damit wird eine nichtlineare Dynamik im Raum der Observablen global
linear repräsentiert, um den Preis unendlicher Dimension. Eine
umfassende moderne Übersicht geben
\cite{bruntonModernKoopmanTheory2022} und
\cite{mauroyKoopmanOperatorSystems2020};
\cite{bruntonModernKoopmanTheory2022} ordnen darin DMD, DMDc, \edmd,
eDMDc und \sindyc{} als Blöcke ein und derselben datengetriebenen
Approximation ein. Diese Sicht stützt den roten Faden dieses
Kapitels.

\cite{kutzDynamicModeDecomposition2016} fassen DMD als
endlichdimensionale Näherung von $\mathcal K$ auf einem
$N$-dimensionalen Zustandsraum zusammen; genauer ist DMD die
datengetriebene Approximation des Koopman-Operators, eingeschränkt auf
lineare Messfunktionen \cite{bruntonModernKoopmanTheory2022}. Formal
werden zwei Snapshot-Matrizen
$\mathbf X_0 = [\mathbf x_1, \dots, \mathbf x_{N-1}]$ und
$\mathbf X_1 = [\mathbf x_2, \dots, \mathbf x_N]$ gebildet; DMD ist dann
genau (\ref{eq.lsq}) mit $\boldsymbol\Theta = \mathbf X_0$ und
$\mathbf Y = \mathbf X_1$:

\begin{equation}
  \label{eq.dmd}
  \mathbf A_{\text{DMD}} \;=\; \mathbf X_1\,\mathbf X_0^+.
\end{equation}



\subsection{Dynamic Mode Decomposition mit Regeleingängen (DMDc)}

Die Eigenwerte von $\mathbf A_{\text{DMD}}$ liefern die
DMD-Eigenfrequenzen, die Eigenvektoren die DMD-Moden.
\cite{proctorDynamicModeDecomposition2016} erweitern DMD um Eingänge zu
DMDc: Die Merkmalsmatrix wird um den Stelleingang ergänzt,
$\boldsymbol\Theta = [\mathbf X_0;\,\mathbf U_0]$, und
$[\mathbf A, \mathbf B]$ werden gemeinsam aus

\begin{equation}
  \label{eq.dmdc}
  [\mathbf A, \mathbf B]
  \;=\;
  \mathbf X_1 \cdot \begin{bmatrix}\mathbf X_0\\ \mathbf U_0\end{bmatrix}^{+}
\end{equation}

geschätzt. DMDc trennt damit Eigendynamik und Stellwirkung
\cite{bruntonModernKoopmanTheory2022}: im vorliegenden Fall die
Trennung von thermischer Eigendynamik und Ventilwirkung. DMDc ist das
\edmd-Pendant zu \sindyc{} und stellt die unmittelbare Verbindung zur
klassischen linearen Zustandsraum\-Systemidentifikation her; die
\emph{Nichtlinearität} bleibt jedoch unerreichbar, solange
$\boldsymbol\Theta$ nur rohe Zustände enthält.


\subsection{Extended Dynamic Mode Decomposition (eDMD) und eDMDc}
\label{ssec.sot.edmd.def}

Genau hier setzt \edmd{} nach \cite{williamsDataDrivenApproximation2015}
an: Statt die Dynamik im niedrigdimensionalen Zustandsraum linear zu
erzwingen, hebt \edmd{} den Zustand über eine Liftung
$\boldsymbol\psi(\mathbf x) = [\psi_1(\mathbf x),\dots,\psi_L(\mathbf x)]^\top$
in einen höherdimensionalen Raum von Observablen, in dem die Dynamik
näherungsweise linear wird: Nichtlinearität wird gegen Dimension
getauscht. Mit $\boldsymbol\Theta = \boldsymbol\psi(\mathbf X_0)$
und $\mathbf Y = \boldsymbol\psi(\mathbf X_1)$ ergibt dieselbe Regression
(\ref{eq.lsq}) eine Matrix
$\mathbf K \in \mathbb R^{L \times L}$ mit

\begin{equation}
  \label{eq.edmd}
  \boldsymbol\psi(\mathbf x_{k+1}) \;\approx\; \mathbf K\, \boldsymbol\psi(\mathbf x_k).
\end{equation}

\edmd{} ist damit dieselbe kleinste-Quadrate-Regression auf einem
gehobenen Zustand und konvergiert im Datenlimit gegen die
Galerkin-Projektion des Koopman-Operators auf den durch
$\boldsymbol\psi$ aufgespannten Unterraum
\cite{williamsDataDrivenApproximation2015, bruntonModernKoopmanTheory2022}.
Zwei Fallstricke sind für diese Arbeit relevant. Ohne einen tatsächlich
Koopman-invarianten Unterraum entstehen Scheineigenwerte, und das
gemeinsame Mitführen von $\mathbf x$ und seiner nichtlinearen Liftung
erzeugt Closure-Probleme
\cite[S.~285--286]{bruntonModernKoopmanTheory2022}.

Um den Kreis zu schließen, wird (analog zum Schritt von DMD zu DMDc)
der Eingang zurückgeholt: Mit
$\boldsymbol\Theta = [\boldsymbol\psi(\mathbf X_0);\,\mathbf U_0]$
entsteht eDMDc mit der gelifteten Update-Gleichung
$\mathbf z_{k+1} = \mathbf A\,\mathbf z_k + \mathbf B\,\mathbf u_k$,
$\mathbf z = \boldsymbol\psi(\mathbf x)$. Der praktische Reiz für die
Regelung: Ein darauf aufgesetztes MPC-Problem ist ein konvexes
quadratisches Programm, dessen Aufwand unabhängig von der Zahl der
Observablen ist \cite[S.~304]{bruntonModernKoopmanTheory2022}.
\cite{kordaLinearPredictorsNonlinear2018} zeigen, dass eDMDc-Modelle für
nichtlineare MPC unter geeigneten Bedingungen asymptotisch konsistent
gegen die wahre Koopman-Repräsentation konvergieren;
\cite{boldDataDrivenMPCStability2025} liefern jüngst
Stabilitätsgarantien für eDMD-basierte MPC. Bemerkenswert ist der von
\cite[S.~305--306]{bruntonModernKoopmanTheory2022} gezeigte
Grenzfall: Ein eDMDc-Modell auf reinen Zeitverzögerungskoordinaten ist
strukturell mit einem ARX-Modell identisch. Damit schließt sich der
Bogen zurück zur Baseline aus Abschnitt~\ref{sec.sot.arx}.





% Bilineare Realisierung. Keys bruderAdvantagesBilinearKoopman2021 und
% yahagiGeneralizedBilinearKoopman2026, beide in Masterarbeit.bib vorhanden.
Zwischen dem in der Stellgröße linearen eDMDc-Prädiktor und einem in
der Stellgröße voll nichtlinearen Modell liegt die bilineare
Realisierung

\begin{equation}
  \label{eq.edmd.bilinear}
  \mathbf z_{k+1} \;=\; \mathbf A\,\mathbf z_k
  + \bigl(\mathbf N\,\mathbf z_k\bigr)\,u_k,
\end{equation}

hier für den in dieser Arbeit vorliegenden skalaren Stelleingang
angeschrieben ($\mathbf N \in \mathbb R^{L\times L}$; bei mehreren
Eingängen tritt je Eingang eine eigene Matrix hinzu). Die Stellwirkung
skaliert mit dem gelifteten Zustand.
\cite{bruderAdvantagesBilinearKoopman2021} zeigen theoretisch und
experimentell, dass bilineare Koopman-Modelle lineare bei gleichem
Datenbudget übertreffen, sobald die Dynamik eingangsaffin ist, die
Stellgröße also über einen zustandsabhängigen Faktor eingeht.
\cite{yahagiGeneralizedBilinearKoopman2026} bauen die bilineare
Realisierung aus reinen Ein-/Ausgangsdaten auf und erproben sie als
Mehrschritt-Prädiktor an einem realen Industriesystem. Für die
Ladeluftstrecke ist diese Form einschlägig, weil die Ventilwirkung mit
der anliegenden Temperaturspanne skaliert
(Gl.~(\ref{eq.grundlagen.gain.epsntu})).
Der Preis ist der
Verlust des konvexen QP: Auf bilinearen Prädiktoren wird das
MPC-Problem sequentiell quadratisch gelöst
\cite{kelmanBilinearModelPredictive2011}.


Die Wahl der Liftung $\boldsymbol\psi$ in (\ref{eq.edmd}) legt fest, welche
Nichtlinearität überhaupt darstellbar ist. Zwei Familien sind für die
vorliegende Strecke einschlägig. Polynomiale Observablen,
$\psi_i(\mathbf x) \in \{1, x_1, x_2, x_1 x_2, x_1^2, \dots\}$, halten die
Liftung konsistent zur Polynombibliothek in \sindy{} und sind die
Standardwahl der verbreiteten Implementierungen.
Zeitverzögerungskoordinaten, $\psi_i(\mathbf x_k) = \mathbf x_{k-i}$, sind
eng mit der klassischen NARX-Featurewahl verwandt und stellen die oben
gezeigte Brücke zur linearen Systemidentifikation her; theoretisch getragen
werden sie vom Satz von Takens, nach dem die Zeitverzögerungseinbettung
eines einzelnen Skalarsignals unter milden Bedingungen eine
differenzierbare Einbettung der Trajektorie liefert.
\cite{bruntonModernKoopmanTheory2022} berichten, dass Delay-Koordinaten für
Regelungszwecke den reinen Monom-Observablen überlegen sind und \edmd{} ohne
Delays trotz Liftung versagen kann.

%%%%%%%%%%%%%%%%%%%%%%%%%%%%%%%%%%%%%%%%%%%%%%%%%%%%%%%%%%%%%%%%%%%%%
%%%%%%%%%%%%%%%%%%%%%%%%%%%%%%%%%%%%%%%%%%%%%%%%%%%%%%%%%%%%%%%%%%%%%
%%  Sparse Identification of Nonlinear Dynamics
%%%%%%%%%%%%%%%%%%%%%%%%%%%%%%%%%%%%%%%%%%%%%%%%%%%%%%%%%%%%%%%%%%%%%
%%%%%%%%%%%%%%%%%%%%%%%%%%%%%%%%%%%%%%%%%%%%%%%%%%%%%%%%%%%%%%%%%%%%%
\section{Sparse Identification of Nonlinear Dynamics (SINDy)}
\label{sec.sot.sindy}

%\subsection{Grundalgorithmus und Optimierungsproblem}
\label{ssec.sot.sindy.basis}

\sindy{} wurde von \cite{bruntonDiscoveringGoverningEquations2016}
eingeführt. Ausgangspunkt ist die Annahme, dass sich die Dynamik
$\dot{\mathbf x} = f(\mathbf x)$ eines realen Systems in einer geeigneten
Bibliothek
$\boldsymbol\Theta(\mathbf x) = [\theta_1(\mathbf x),\dots,\theta_p(\mathbf x)]$
nichtlinearer Kandidatenfunktionen durch wenige dominante Terme
darstellen lässt:

\begin{equation}
  \label{eq.sindy.def}
  \dot{\mathbf X} \;=\; \boldsymbol\Theta(\mathbf X)\,\boldsymbol\Xi,
  \qquad
  \boldsymbol\Xi \in \mathbb{R}^{p \times n},\quad \|\boldsymbol\Xi\|_0 \ll p\,n.
\end{equation}

Im Vergleich zur Klammer (\ref{eq.lsq}) sind zwei Dinge neu: Die
Zielgröße ist die Zeitableitung $\dot{\mathbf X}$ (nicht der
Folgezustand), und die Lösung wird dünn besetzt erzwungen. Die
$\ell_0$-Pseudonorm $\|\boldsymbol\Xi\|_0$ zählt dabei die von null
verschiedenen Einträge; gefordert sind deutlich weniger aktive Terme
als die $p\,n$ möglichen Einträge, die $p$ Kandidatenfunktionen für
$n$ Zustandsgleichungen aufspannen. Die Schätzung
von $\boldsymbol\Xi$ erfolgt durch eines der folgenden Verfahren:

\begin{itemize}
  \item \textbf{Sequentially thresholded least squares (STLSQ).}
    Alterniert zwischen unbeschränkter Least-Squares-Lösung und
    Schwellwert-Setzen aller Koeffizienten $|\xi_i| < \lambda$ auf Null
    \cite[Algorithmus~1]{bruntonDiscoveringGoverningEquations2016}.
    Einfach, schnell, mit einem skalaren Hyperparameter.
  \item \textbf{LASSO} ($\ell_1$-regularisierte Regression): löst
    $\min_{\boldsymbol\xi} \|\dot{\mathbf X} - \boldsymbol\Theta\boldsymbol\xi\|_2^2
    + \lambda\|\boldsymbol\xi\|_1$ und liefert eine weichere
    Sparsifizierung; Elastic Net kombiniert $\ell_1$ und $\ell_2$ und ist
    bei stark korrelierten Bibliotheksspalten robuster.
\end{itemize}

%\subsection{Wahl der Bibliothek}
\label{ssec.sot.sindy.bib}

Die Bibliothek $\boldsymbol\Theta$ bestimmt, welche Strukturen sich
identifizieren lassen. Übliche Bausteine sind Polynome bis Grad $d$,
trigonometrische Funktionen, rationale Terme und Indikatoren. Hinzu treten
stückweise lineare Terme (Hinge), die eine Kennlinie mit Knick abbilden,
und der Zustand einer exponentiellen Glättung, mit dem ein träger Anteil
eines Eingangs als eigene Spalte geführt wird.
\cite[Kap.~2]{bruntonDataDrivenScience2017} und
\cite{desilvaPySINDyPythonPackage2020} geben eine systematische Übersicht.
Für physikalisch motivierte Probleme empfehlen
\cite{loiseauConstrainedSparseGalerkin2018}, die Bibliothek aus den
Erhaltungsgleichungen abzuleiten und harte Gleichungsnebenbedingungen
einzubauen. Diesen Weg greift die vorliegende Arbeit auf; er sichert
zugleich die Dimensionstreue der Kandidatenterme, die sonst eine
vorgeschaltete Buckingham-$\pi$-Analyse herstellen müsste
\cite{wulffPhysicsinformedRegressionModels2025}.


Die physik-informierte Linie reicht über die Bibliothekswahl hinaus. In
Grey-Box-Strategien dient das physikalische Modell als Prior, während der
gelernte Anteil die unmodellierten Residuen trägt
\cite{subramanianWhiteboxMachineLearning2021,
ellisMachineLearningEnhanced2021}. Eine online-aktualisierbare Variante
liegt mit OASIS-P vor \cite{bhadrirajuOASISPOperableAdaptive2021}. Die
Nachführung eines Modells über Verschmutzung und Sensordrift hinweg ist
nicht Gegenstand dieser Arbeit.


\subsection{SINDy mit Regeleingängen (SINDyc)}
\label{ssec.sot.sindyc}

\sindyc{} erweitert die Formulierung um Eingänge $\mathbf u$:

\begin{equation}
  \label{eq.sindyc.def}
  \dot{\mathbf X} \;=\; \boldsymbol\Theta(\mathbf X, \mathbf U)\,\boldsymbol\Xi.
\end{equation}

Die Methode geht zurück auf
\cite{bruntonSparseIdentificationNonlinear2016}; ein gut lesbares
Tutorial gibt \cite{faselSINDyControlTutorial2021}.
Bild~\ref{fig.sindyc.pipeline} zeigt den methodischen Ablauf im
Überblick: Aus den gemessenen Zeitreihen wird eine nichtlineare
Bibliothek aufgebaut, deren Koeffizienten anschließend per sparser
Regression bestimmt werden. Wesentlich für diese Arbeit: \sindyc{} setzt
voraus, dass die Eingänge tatsächlich Stellgrößen sind, die das System
exzitieren. Ob der Charge-Air-Pfad diese Voraussetzung erfüllt, ist eine
Frage an den Datenbestand (vgl.\
\cite{boddupalliKoopmanOperatorsGeneralized2019} und
Kapitel~\ref{ch.datenbasis}); beantwortet wird sie in
Abschnitt~\ref{sec.erg.voranalyse} und in
Abschnitt~\ref{sec.erg.clarbiter}.

\begin{figure}[bht]
  \begin{center}
    \includegraphics[width=\textwidth]{sindyc_pipeline}
    \caption{Schematische Darstellung der \sindyc-Pipeline. Aus den
      gemessenen Zeitreihen von Zustand $\mathbf X$, Stellgröße
      $\mathbf U$ und Störgrößen $\mathbf D$ wird eine Bibliothek
      $\boldsymbol\Theta(\mathbf X, \mathbf U, \mathbf D)$ nichtlinearer
      Kandidatenfunktionen aufgebaut. Die anschließende sparse Regression
      $\min_{\boldsymbol\xi_k}\,\|\dot{\mathbf x}_k -
      \boldsymbol\Theta^{\!\top}(\mathbf X, \mathbf U, \mathbf D)\,
      \boldsymbol\xi_k\|_2 + \lambda\|\boldsymbol\xi_k\|_1$ liefert die
      dünn besetzte Koeffizientenmatrix $\boldsymbol\Xi$ und damit das
      identifizierte System.}
    \label{fig.sindyc.pipeline}
  \end{center}
\end{figure}


\subsection{Esemble-Erweiterungen}
\label{ssec.sot.sindy.robust}

Die in \cite{bruntonDiscoveringGoverningEquations2016} präsentierte
Standardform ist gegenüber Messrauschen und kurzen Trainingsfenstern
empfindlich. Drei Erweiterungen sind in der Literatur konsolidiert.
Ensemble-\sindy{} \cite{faselEnsembleSINDyRobustSparse2022} aggregiert
Modelle aus Bootstrap-Stichproben der Bibliothek und der Datenmatrix und
liefert Inklusionswahrscheinlichkeiten je Kandidatenterm. Weak-\sindy{}
\cite{messengerWeakSINDyPartial2021} ersetzt die punktweise Ableitung durch
eine schwache Form, in der die Gleichung mit einer Testfunktion gewichtet
und über ein Zeitfenster integriert wird, so dass die Differentiation
analytisch auf die Testfunktion fällt und die Verstärkung hochfrequenten
Rauschens entfällt. Die Modellselektion über
Informationskriterien \cite{dongImprovedSparseIdentification2023,
manganModelSelectionDynamical2017} ersetzt den Schwellwert in STLSQ durch
AIC/BIC oder eine Pareto-Front aus Trainingsfehler und Sparsität; sie tritt
hier nicht als Wahlregel auf, sondern allein als nachgelagerte
Sparsamkeitsprüfung. Woran zwei Fassungen unterschieden
werden, legt die Bewertungsleiter in Abschnitt~\ref{ssec.meth.gates}
fest.


\label{ssec.sot.sindy.discrete}

Gl.~(\ref{eq.sindy.def}) ist in der Brunton-Originalform kontinuierlich
angeschrieben: $\dot{\mathbf x} = \boldsymbol\Theta(\mathbf x)\boldsymbol\Xi$.
In der praktischen Implementierung steht aber eine zeitdiskrete
Zustandsfolge $\{\mathbf x_k\}$ zur Verfügung. Zwei Formulierungen sind
verbreitet:

\begin{itemize}
  \item \textbf{Kontinuierlich (Standardform).} Schätze
    $\dot{\mathbf x}_k$ numerisch und identifiziere $\boldsymbol\Xi$ in
    der kontinuierlichen Form. Vorteil: physikalisch lesbares Modell mit
    Größen wie $[\mathrm{K}/\mathrm{s}]$. Nachteil: numerische
    Differentiationsempfindlichkeit gegen Rauschen.
  \item \textbf{Diskret (map-Form).} Identifiziert wird die Update-Abbildung
    $\mathbf x_{k+1} = \mathbf x_k + \Delta t\,
    \boldsymbol\Theta(\mathbf x_k)\boldsymbol\Xi$ direkt. Vorteil: keine
    Ableitung nötig. Nachteil: implizite Abhängigkeit von $\Delta t$
    macht das Modell weniger universell.
\end{itemize}

\begin{figure}[htbp]
  \centering
  \includegraphics[width=\textwidth]{kapitel3/pictures/Esemble_Methoden.png}
  \caption[Ensemble-SINDy: Daten- und Bibliotheksstichproben]{Ensemble-SINDy:
    Daten-Bagging und Library-Bagging. Übernommene Darstellung aus
    \cite[Abb.~2]{faselEnsembleSINDyRobustSparse2022};
    $\mathbf U$ als Zustandsdatenmatrix in der Notation der Quelle,
    nicht als Stellgrößenmatrix dieser Arbeit.}
  \label{fig.sot.ensemble}
\end{figure}




%%%%%%%%%%%%%%%%%%%%%%%%%%%%%%%%%%%%%%%%%%%%%%%%%%%%%%%%%%%%%%%%%%%%%
%%%%%%%%%%%%%%%%%%%%%%%%%%%%%%%%%%%%%%%%%%%%%%%%%%%%%%%%%%%%%%%%%%%%%
%%  SINDy vs eDMD
%%%%%%%%%%%%%%%%%%%%%%%%%%%%%%%%%%%%%%%%%%%%%%%%%%%%%%%%%%%%%%%%%%%%%
%%%%%%%%%%%%%%%%%%%%%%%%%%%%%%%%%%%%%%%%%%%%%%%%%%%%%%%%%%%%%%%%%%%%%

\section{Gemeinsamkeiten und Grenzen der Verfahren}
\label{ssec.sot.edmd.vs.sindy}

Die bisher getrennt behandelten Verfahren lösen zu einem großen Teil
dasselbe Problem. \arx, \sindyc{} und \edmdc{} sind linear in ihren
Parametern und bestimmen eine Koeffizientenmatrix nach (\ref{eq.lsq}).
Sie unterscheiden sich in drei Entwurfsentscheidungen: der Bibliothek
$\boldsymbol\Theta$, dem Lernziel $\mathbf Y$ und der Regularisierung,
mit der die Koeffizienten beschränkt oder ausgewählt werden.

\sindy{} sucht eine sparsame Darstellung der Zeitableitung im Zustand,
\edmd{} eine lineare Darstellung der Folgebeobachtung im gehobenen Raum
((\ref{eq.sindy.def}) und (\ref{eq.edmd})). Der Unterschied liegt im
Lernziel und in der Bibliothek, nicht im Schätzverfahren.

Für die beiden oberen Verfahrensfamilien kommt zu diesen Entscheidungen
eine Reihe gemeinsamer Grenzen hinzu, die aus der Wahl der Bibliothek und
aus der Schätzung über kleinste Quadrate folgen. Sie stehen in
Abschnitt~\ref{sec.sot.grenzen}.

\label{sec.sot.grenzen}

Eine ausgewogene Würdigung beider Verfahrensfamilien muss auch ihre
Grenzen benennen:

\begin{itemize}
  \item \textbf{Empfindlichkeit gegen Messrauschen.} Insbesondere
    \sindy{} mit numerischer Ableitung reagiert empfindlich; Ensemble-
    und Weak-Varianten mildern, lösen das Problem aber nicht vollständig
    \cite{faselEnsembleSINDyRobustSparse2022, messengerWeakSINDyPartial2021}.
  \item \textbf{Abhängigkeit von der Bibliothek bzw.\ Liftung.} Bei
    unzureichender Bibliothek wird kein konsistentes Modell
    identifiziert. Eine systematische Bibliothekswahl ist nach wie vor
    offener Forschungspunkt \cite{kaiserSparseIdentificationNonlinear2018}.
  \item \textbf{Extrapolationsfähigkeit.} Da die Bibliothek lokal aus
    Daten gefittet wird, ist Extrapolation in nicht beprobte
    Betriebsbereiche nur eingeschränkt zuverlässig
    \cite{atamAdvancedAirPath2018}.
  \item \textbf{Stabilitätsgarantien.} Für \edmd-MPC liegen jüngst
    formale Garantien vor \cite{boldDataDrivenMPCStability2025}; für
    \sindy-MPC ist die Theorie weniger ausgereift.
\end{itemize}

Diese Grenzen begründen, warum die vorliegende Arbeit beide Familien
parallel prüft und sich nicht vorab auf eine Methodenwahl festlegt.


\section{Closed-Loop-Identifikation}
\label{closed_loop_identifkation}

Im Feld folgt die Ventilstellung der Last. Jeder verfügbare Feldmesswert
entsteht deshalb unter Rückführung, in der Stellgröße und Störung über
den Regler korreliert sind.


textbf{Identifizierbarkeit unter Closed-Loop-Bedingungen.}
    \cite{dahdahClosedloopKoopmanOperator2024} zeigen, dass naives
    \edmd{} bei rückgekoppelten Eingängen verzerrte Operatoren liefert;
    eine analoge Sensitivität ist für \sindy{} zu erwarten.


Das Problem ist in den Erstveröffentlichungen der hier verwendeten
Verfahren selbst benannt.
Für \sindyc{} halten
\cite[S.~713]{bruntonSparseIdentificationNonlinear2016} fest: Wird die
Stellgröße per Rückführung eingestellt, korreliert sie mit dem Zustand,
die Wirkung der Stellgröße ist von der Eigendynamik nicht mehr zu
unterscheiden, und die Regression wird schlecht konditioniert. Für DMDc
beschreiben \cite[S.~153]{proctorDynamicModeDecomposition2016} denselben
Befund: Bei zustandsabhängiger Rückführung $u_k = K\,\mathbf x_k$ ist
$\mathbf B$ mit dem dort beschriebenen Verfahren nicht rekonstruierbar;
in Gl.~(\ref{eq.dmdc}) verliert die Merkmalsmatrix
$[\mathbf X_0;\,\mathbf U_0]$ dann ihren Zeilenrang. Ein aufaddiertes
mittelwertfreies Störsignal, $u_k = K\,\mathbf x_k + \delta_k$, genügt:
Es bricht die Symmetrie der Rückführung und macht $\mathbf A$ und
$\mathbf B$ getrennt schätzbar. Ein solches Störsignal setzt allerdings
einen Eingriff in die Anlage voraus.
% Trennung SdT/Methodik (Copilot 02.09.2026, Anweisung Thomas): der Verweis
% auf den eigenen Diesel/Gas-Modussprung stand hier und ist nach Kap. 5
% verschoben. Stand der Technik traegt nur das Allgemeine.
% Bildvorschlag (Blockschaltbild des geschlossenen Kreises, VISIO) steht in
% notizen\CH3_SDT_NUTZUNG_0309.md § 4.


Die Verfahren zur Identifikation unter Rückführung lassen sich in drei
Schulen ordnen \cite{forssellClosedloopIdenticationRevisited1999}. Die
direkte Methode schätzt aus Stellgröße und Ausgang, als wäre der Kreis
offen. Konsistent ist sie nur, wenn das wahre Störmodell $H_0$ im
Modellsatz liegt \cite{gustavssonIdentificationProcessesClosed1977}.
% Quelle: karimiComparisonClosedloopIdentification1998, S. 161
% ALT (Copilot 01.09.): "zur Differenz $H-\hat H$" — $H$ war nicht eingeführt,
% oben heißt das wahre Störmodell $H_0$.
Liegt $H_0$ außerhalb, bleibt die Verzerrung der Strecke proportional zur
Rauschvarianz und zur Differenz $H_0-\hat H$ zwischen wahrem und
angesetztem Störmodell
\cite{karimiComparisonClosedloopIdentification1998}. Davon unabhängig
wirkt eine zweite Ursache, die fehlende Anregung.
% Quelle: lovlandClosedloopIdentificationChallenge2025, Theorem 1, Gl. (26)
Geht der Anteil der Sollwertvarianz an der Gesamtvarianz gegen null,
wandert der Grenzwert des Schätzers gegen die negativ invertierte
Reglerverstärkung $-K^{-1}$
\cite{lovlandClosedloopIdentificationChallenge2025}.
% REPARATUR (Copilot, 02.09.2026): der Satzanfang "Ein ARX-Praediktor" war in die
% Kommentarzeile geraten, der gesetzte Satz begann deshalb mit "hoher Ordnung".
Ein ARX-Prädiktor
hoher Ordnung (HO-ARX) oder eine Box-Jenkins-Struktur ist die übliche
Abhilfe gegen die erste Ursache \cite{ljungSystemIdentificationTheory1999}.

Die indirekte Methode identifiziert zuerst den geschlossenen Kreis vom
Sollwert zum Ausgang und rechnet über den bekannten Regler auf die
Strecke zurück \cite{forssellClosedloopIdenticationRevisited1999}. Sie
verlangt einen linearen und exakt bekannten Regler
\cite{agueroVirtualClosedLoop2011}. Ohne bewegten Sollwert fehlt der
Rückrechnung die Anregung. Joint-Input-Output-Verfahren und die
dual-Youla-Parametrierung verallgemeinern beide Wege über einen
virtuellen Regler, mit dem sich der Kompromiss zwischen der Genauigkeit
des Rauschmodells und der Reglerkenntnis parametrieren lässt
\cite{agueroVirtualClosedLoop2011}%
\anno{Quelle nachtragen: Hansen 1989, Erstquelle der
dual-Youla-Parametrierung}. Als Subspace-Variante für den geschlossenen
Kreis tritt die prädiktorbasierte Subspace-Identifikation hinzu (PBSID,
engl.\ Predictor-Based Subspace Identification)%
\anno{Quelle nachtragen: Chiuso 2007, Erstdefinition PBSID}. Für
rückgekoppelt stabilisierte Systeme liegt zudem eine neuronale Ausprägung
vor \cite{boroujeniNeuralIdentificationFeedbackStabilized2025}.

Die Instrumentalvariablen-Schätzung und das Zwei-Stufen-Verfahren nutzen
beide den Sollwert als exogene Größe, die mit dem Ausgangsrauschen
unkorreliert ist \cite{gilsonInstrumentalVariableMethods2005}. Das
Zwei-Stufen-Verfahren projiziert die Stellgröße zuerst auf den Sollwert
und schätzt die Strecke aus dem projizierten Anteil
\cite{vandenhofIdentificationDynamicModels2013}%
\anno{Quelle nachtragen: Schrama 1993, Erstquelle des
Zwei-Stufen-Verfahrens}. Die Einordnung fällt je nach Quelle verschieden
aus. Das Zwei-Stufen-Verfahren gilt einmal als indirektes Verfahren
\cite{karimiComparisonClosedloopIdentification1998} und einmal als
Ausprägung der Joint-Input-Output-Familie
\cite{forssellClosedloopIdenticationRevisited1999}, die
Instrumentalvariable wird als eigene Klasse geführt
\cite{gilsonInstrumentalVariableMethods2005}.
% Quelle: karimiComparisonClosedloopIdentification1998, S. 160
Die Verzerrung verteilt sich über die Frequenz, und ihre Gewichtung hängt
vom Regler ab \cite{karimiComparisonClosedloopIdentification1998}. Bei
Instrumentalvariablen- und Zwei-Stufen-Verfahren bestimmen die Stärke des
Instruments und die gewählte Modellordnung, wie viel Verzerrung übrig
bleibt \cite{gilsonInstrumentalVariableMethods2005}. Einen dritten Weg
gehen Verfahren zur Korrektur der Verzerrung, die den Verzerrungsterm
mitschätzen
\cite{mejariBiascorrectionMethodClosedloop2018,mandloiMethodsClosedLoop2015};
für den geschlossenen Kreis ist die Äquivalenz der bias-eliminierten
Kleinste-Quadrate-Schätzung (BELS) und der Instrumentalvariablen-Schätzung
gezeigt \cite{gilsonRelationBiaseliminatedLeastsquares2001}.



% Quelle: bittencourtAlgorithmFindingProcess2015, S. 377 (1,45 % markierte Abtastwerte), S. 369 (Ausregelbetrieb per Konstruktion ausgeschlossen), S. 378 (Ausschlussgrund im geschlossenen Kreis)
Wie selten die Bedingung im Anlagenbetrieb erfüllt ist, zeigt eine
Auswertung von 195 Regelkreisen über 37 Monate. Als potenziell brauchbar
markiert der Suchalgorithmus 1,45\,\% der Abtastwerte, wobei
Ausregelbetrieb bei konstantem Sollwert per Konstruktion nicht durchsucht
wird \cite{bittencourtAlgorithmFindingProcess2015}. Die Zahl misst damit
die Ausbeute dieses Verfahrens, nicht den Informationsgehalt von
Routinedaten schlechthin. Im geschlossenen Kreis ist der fehlende
Sollwertsprung der häufigste Ausschlussgrund.
% Quelle: bittencourtAlgorithmFindingProcess2015, Tabelle 2 (S. 378); die 40 min je Kreis sind eigene Rechnung aus dieser Tabelle, siehe EHRGQL83_bittencourt2015_*.md Z. 61, 77 und 102
Für den hier einschlägigen Kreistyp fällt der Befund noch enger aus. Auf
die 50 Temperaturkreise entfallen über die 37 Monate rund 40 min
verwertbare Daten je Kreis, gerechnet aus der dortigen Tabelle 2.

Das Verwerfen schwach informativer Abschnitte senkt den Gesamtfehler,
sofern das Modell in Gleichungsfehlerstruktur mit gemeinsamem
Nennerpolynom vorliegt und die Daten aus dem offenen Kreis stammen
\cite{carretteDiscardingDataMay1996}. Bei untermodelliertem Rauschen
verzerren die stationären Abschnitte die Modellparameter; daraus folgt
kein allgemeines Argument gegen große Datenmengen. Wie eine
Anregung auszulegen ist, die den laufenden Betrieb wenig stört und für
den Reglerentwurf dennoch reicht, behandelt der Entwurf des
kostenminimalen Experiments \cite{bomboisLeastCostlyIdentification2006}.

Auf der Koopman-Seite ist der Befund enger gefasst. Ein rückgekoppelter
Eingang $u = h(\mathbf x)$ macht Eigendynamik und Stellwirkung
ununterscheidbar \cite{proctorGeneralizingKoopmanTheory2018}.
Bei bekanntem Regler lassen sich Kreis und Strecke gemeinsam und unter
Nebenbedingungen schätzen, bewertet am Vorhersagefehler des geschlossenen
Kreises \cite{dahdahClosedloopKoopmanOperator2024}.


%%%%%%%%%%%%%%%%%%%%%%%%%%%%-------------------------------------------
%Identifikation im geschlossenen Regelkreis
%%%%%%%%%%%%%%%%%%%%%%%%%%%-------------------------------------------

%\section{Identifikation im geschlossenen Regelkreis}
%\label{ssec.grundlagen.clid}

In meisten Industriellen Systeme laufen unter einer regelung und die Eingangssignae der Systeme sind teils durch ein Rückgabe anderer signale bestimmt. Sobald dies der Fall ist, liegt ist das Stellsignal korreliert mit der Regelgröße. Dadurch lässt sich schwer der genaue Einfluss der Stellgröße und der REgelgröße unterscheiden da sie sich stark abhängig voneinander verhalten.

Um dies aufzulösen gibt es zwei prinzipielle Vorgehensweisen:

Direkte Motehode:

Man ignoriert es einfach und bestimmt direkt aus den messdaten

Indirkete Methdoe;


Messdaten. Die Rückkopplung bindet die Stellgröße an den Verlauf der
Zielgröße. Die Strecke wird für die folgende Betrachtung als lineares
Modell mit additivem Störanteil geschrieben,
\begin{equation}
  \label{eq.grundlagen.cl.strecke}
  y(t) \;=\; G(q)\,u(t) \;+\; H(q)\,e(t),
\end{equation}
wobei $q$ den Verschiebeoperator bezeichnet, $e(t)$ weißes Rauschen ist
und die Störübertragung $H(q)$ als monisch und stabil vorausgesetzt
wird. Der Regler bildet die Stellgröße aus der Regelabweichung,
\begin{equation}
  \label{eq.grundlagen.cl.regler}
  u(t) \;=\; C(q)\,\bigl(r(t) - y(t)\bigr).
\end{equation}
Einsetzen von Gl.~\eqref{eq.grundlagen.cl.strecke} in
Gl.~\eqref{eq.grundlagen.cl.regler} und Auflösen nach $u(t)$ liefert
\begin{equation}
  \label{eq.grundlagen.cl.u}
  u(t) \;=\; S(q)C(q)\,r(t) \;-\; S(q)C(q)H(q)\,e(t),
  \qquad S(q) = \bigl(1 + C(q)G(q)\bigr)^{-1}.
\end{equation}
Über den zweiten Term führt die Stellgröße die Störung des Prozesses
mit sich, weil der Regler auf sie antwortet.
% ===== ERGÄNZT Copilot 01.09.2026 (Fassung n. Lektor) =====
$S(q)$ ist die Empfindlichkeitsfunktion des geschlossenen Kreises. Ihr
Frequenzgang $|S(\mathrm e^{\mathrm j\omega})|$ gibt an, ob der Kreis
eine Störung unterdrückt ($|S| < 1$), unverändert durchlässt
($|S| \approx 1$) oder verstärkt ($|S| > 1$). Die Auswertung in
Abschnitt~\ref{sec.erg.clarbiter} trägt die Größe über der
Periodendauer auf und greift auf diese Deutung zurück.

Der Regressor $\varphi(t)$ eines in den Parametern linearen Modells
enthält verzögerte Werte von $y$ und $u$, der Vektor $\theta_0$ die
wahren Parameter. Die Schätzung über kleinste Quadrate strebt mit
wachsender Datenlänge gegen
\begin{equation}
  \label{eq.grundlagen.cl.bias}
  \hat{\theta} \;\longrightarrow\; \theta_0 \;+\;
  \bigl[\bar{\mathrm{E}}\{\varphi\varphi^{\top}\}\bigr]^{-1}\,
  \bar{\mathrm{E}}\{\varphi\,v\}, \qquad v(t) = A(q)H(q)e(t).
\end{equation}
Dabei ist $v(t)$ der Gleichungsfehler des Modells mit Nennerpolynom
$A(q)$, nicht die Störung $H(q)e(t)$ selbst.
% Quelle: Ljung 1999, System Identification, 2. Aufl., Kap. 8 (Konvergenz und Konsistenz der LS-Schaetzung)
Der Kreuzterm $\bar{\mathrm{E}}\{\varphi v\}$ verschwindet im offenen
Kreis nur, wenn $v(t)$ weiß ist oder der Regressor keine verzögerten
Ausgänge enthält; andernfalls ist auch die offene Schätzung verzerrt
\cite[Kap.~8]{ljungSystemIdentificationTheory1999}. Im geschlossenen
Kreis ist er nach Gl.~\eqref{eq.grundlagen.cl.u} von null verschieden.
% Quelle: Forssell und Ljung 1999, Automatica 35(7), S. 1215 bis 1241 (Bias der direkten Methode, Kreuzspektrum); Einzelseite im .bib nicht belegt
Im Frequenzbereich steht anstelle des Kreuzterms das Kreuzspektrum
$\Phi_{ue}(\omega)$, das allein ohne Rückführung verschwindet
\cite{forssellClosedloopIdenticationRevisited1999}. Konsistenz der
direkten Methode bleibt erreichbar, verlangt aber ein Rauschmodell, das
die wahre Störübertragung $H_0(q)$ enthält
% Quelle: Gustavsson, Ljung und Söderström 1977, Automatica 13(1), S. 59 bis 75 und Ljung 1999
\cite{gustavssonIdentificationProcessesClosed1977,%
ljungSystemIdentificationTheory1999}.

Beide Terme von Gl.~\eqref{eq.grundlagen.cl.bias} stehen für je eine
Fehlerquelle. Ein singuläres oder schlecht konditioniertes
$\bar{\mathrm{E}}\{\varphi\varphi^{\top}\}$ kennzeichnet den Anregungs-
und Vorsteuerungsfall, ein von null verschiedener Kreuzterm den
Rückkopplungsfall. Die Aktuatorsättigung wirkt auf beide Terme
gegenläufig: sie beseitigt den Kreuzterm und entartet zugleich die
Anregungsmatrix. Die Gleichung gilt für jedes in den Parametern lineare
Modell und damit auch für \sindyc und \edmdc. Bei \sindyc setzt die
Schwellenwertauswahl der Terme auf der verzerrten Lösung auf, sodass die
Verzerrung dort zusätzlich die Termauswahl trifft.

Ob ein Datensatz die gesuchten Parameter festlegt, entscheidet die
Anregung. Ein quasistationäres Signal heißt persistent anregend der
Ordnung $n$, wenn die aus $r_u(\tau) = \bar{\mathrm{E}}\{u(t)u(t-\tau)\}$
ohne Mittelwertabzug gebildete Matrix
\begin{equation}
  \label{eq.grundlagen.cl.pe}
  \bar{R}_n(u) \;=\;
  \begin{bmatrix}
    r_u(0)   & \cdots & r_u(n-1) \\
    \vdots   & \ddots & \vdots   \\
    r_u(n-1) & \cdots & r_u(0)
  \end{bmatrix}
\end{equation}
% Quelle: Ljung 1999, System Identification, 2. Aufl., Kap. 13 (Experimententwurf, Informativitaet, persistente Anregung)
positiv definit ist \cite[Kap.~13]{ljungSystemIdentificationTheory1999}.
% Quelle: Gevers, Bazanella, Bombois und Miskovic 2009, IEEE TAC 54(12), S. 2828 bis 2840
Dieselbe Frage ist als Rangproblem der Informationsmatrix formuliert und
liefert so ein Kriterium für ausreichende Anregung
\cite{geversIdentificationInformationMatrix2009}. Last und Störung
wirken ebenfalls exogen. Stellbar und von der Störung unabhängig ist
allein der Sollwert. Bleibt er unverändert, verliert die
Informationsmatrix ihren Rang in der Richtung des Stellkanals, und der
Fehler wird mit wachsender Datenmenge nicht kleiner.
% Quelle: Ljung 1999, System Identification, 2. Aufl., Abschnitt 13.4 (Identifiability under closed-loop operation)
Ein Regler hinreichend hoher Ordnung oder ein Regler, der zwischen
mehreren Gesetzen umschaltet, macht die Daten auch bei festem Sollwert
informativ \cite[Abschn.~13.4]{ljungSystemIdentificationTheory1999}.
Darauf stützt sich der Modussprung als natürliches Instrument in
Abschnitt~\ref{sec.erg.feld}. Informativität sichert die
Identifizierbarkeit, nicht die Konsistenz eines Schätzers; die Bedingung
an das Störmodell aus Gl.~\eqref{eq.grundlagen.cl.bias} bleibt bestehen.

In dieser Arbeit treten zwei Sonderfälle auf. Im ersten stellt
eine Vorsteuerung die Stellgröße aus einer gemessenen Störgröße,
näherungsweise $\valveCA \approx f(\text{Last})$. Diese Störgröße treibt
zugleich die Strecke, weshalb dasselbe Rangargument greift wie bei einer
Rückführung über den Zustand
% Quelle: Proctor, Brunton und Kutz 2018, SIAM J. Appl. Dyn. Syst. 17(1), S. 909 bis 930
\cite{proctorGeneralizingKoopmanTheory2018}. Stellgröße und Regressor
werden kollinear, und eine größere Datenmenge ändert daran nichts.
% ===== ERGÄNZT Copilot 01.09.2026 (Fassung n. Lektor) =====
Diese Verwechselbarkeit von Stell- und Störwirkung wird im Folgenden
als Konfundierung bezeichnet; ihr Grad wird in
Kapitel~\ref{ch.methodik} als Bestimmtheitsmaß der Stellgröße gegen
die Last gemessen.

Der zweite Sonderfall ist die Sättigung des Aktuators. Steht das Ventil
am Anschlag, wirkt keine Rückführung mehr auf die Stellgröße, und der
Kreuzterm verschwindet. Der erste Term entartet dafür, denn eine
konstante Stellgröße ist persistent anregend der Ordnung 1 und liefert
allein den stationären Anteil der Stellwirkung. Dieser Anteil ist vom
Absolutglied des Modells nicht zu trennen. Der Störpfad bleibt aus
solchen Abschnitten identifizierbar, der Stellkanal nicht.

Eine weitere Bedingung betrifft die Zeitskala. Arbeitet der Regler
schneller als der Abtasttakt des Modells, liegt zwischen zwei
Abtastzeitpunkten bereits eine vollständige Rückkopplung. Deren Wirkung
faltet sich in den Stellkanal ein, ohne dass die Modellstruktur einen
Platz dafür vorsieht. Die Modellstruktur ist in diesem Fall
fehlspezifiziert, vergleichbar einem Aliasing der Rückführung. Wie stark
dieser Anteil wiegt, zeigt Abschnitt~\ref{sec.erg.clarbiter}.

% ===== NEUFASSUNG Copilot 01.09.2026 (Dreiteilung; Altfassung entfernt
% ===== 03.09., s. ch2.tex.bak_pre_4d0309) =====
Ein Streckenmodell hat hier drei getrennt prüfbare Eigenschaften: das
Niveau (den stationären Wert, auf den die Prognose einläuft), die
Dynamik (Zeitkonstante und Totzeit) und die Stellwirkung (die
Verstärkung des Ventilkanals $\valveCA$). Der beschriebene
Rückführmechanismus betrifft die beiden letzten; sie hängen daran, ob
die Stellgröße im Datensatz unabhängig von der Störung variiert. Das
Niveau gehört nicht dazu: Es entsteht als Arbeitspunkt-Offset aus
anderer Ursache und bleibt unabhängig davon stellbar. In der Sättigung
fallen Niveau und stationärer Stellanteil allerdings zusammen (siehe
oben); dort trägt die Trennung nicht. Auf diese Dreiteilung greifen
die Bewertungsleiter (Abschnitt~\ref{ssec.meth.gates}), die
Modellbewertung (Kapitel~\ref{ch.ergebnisse}) und der
Störgrößenbeobachter des Reglers (Abschnitt~\ref{sec.methodik.mpc})
zurück. Welche dieser Eigenschaften sich aus Regelkreisdaten
zurückgewinnen lässt, prüft Abschnitt~\ref{sec.erg.clarbiter} gegen
die am Prüfstand bekannte Wahrheit.
Abschnitt~\ref{sec.erg.feld} überträgt die Frage auf den Feldbetrieb.
















%%%%%%%%%%%%%%%%%%%%%%%%%%%%%%%%%%%%%%%%%%%%%%%%%%%%%%%%%%%%%%%%%%%%%
%%  Diskusion im Bezug zum Thema
%%%%%%%%%%%%%%%%%%%%%%%%%%%%%%%%%%%%%%%%%%%%%%%%%%%%%%%%%%%%%%%%%%%%%

\section{Vorhandene Arbeiten mit Bezug zum Thema}
\label{sec.sot.application}
Intern wurde bereits eine Arbeit zu der Modellierung der Ansteuerung des Dreiwegeventils geschrieben. Dort wird die LT Cooling Water Temperatur als konstant angenommen. \cite{rupprechtAnalysisSimulationOptimisation2016}



Datengetriebene Modellierung und Regelung des Luftpfades wurde primär für
PKW- und Nutzfahrzeug\-Dieselmotoren untersucht. Wichtige Beiträge:

\begin{itemize}
  \item \cite{atamAdvancedAirPath2018} entwickelt einen erweiterten
    LPV-Regler für den Diesel-Luftpfad und überführt nichtlineare
    Modelle in eine parameterabhängige lineare Form.
  \item \cite{aranDieselEngineAirpath2020, aranFlexibleRobustControl2019}
    kombinieren datengetriebene Störbeobachter mit Sliding-Mode- und
    GPR-basierten Reglern für serienreife Steuergeräte. Die Dissertation
    \cite{aranFlexibleRobustControl2019} ist methodisch die direkteste
    Referenzarbeit zur datengetriebenen Luftpfad-Modellierung eines
    schwergängigen Diesel-Antriebs; sie verwendet jedoch
    Multi-Input-Single-Output-NFIR-Modelle für MAF und MAP, nicht
    \sindy{} oder \edmd, und zielt auf PKW-/Lkw-Größenklassen.
  \item \cite{bardosControlOrientedAir2013} und
    \cite{moriyasuDieselEngineAir2019} stellen regelungsorientierte
    Modelle für aufgeladene Dieselmotoren vor.
  \item Für \emph{Großmotoren} bzw.\ Schiffsdiesel sind die Arbeiten von
    \cite{baldiDevelopmentCombinedMean2015} (kombiniertes MVEM/0D-Modell
    für einen Viertakt-Schiffsdiesel) und
    \cite{theotokatosCycleMeanValue2010} maßgeblich. Beide Beiträge sind
    White-Box-getrieben und liefern damit den Kontrast zur
    datengetriebenen Sicht dieser Arbeit.
  \item \cite{beranModelbasedApproachControl2021} ist die für die
    Ladelufttemperatur direkteste Arbeit: eine Regelstrategie für ein
    Ladeluftkühlkonzept mit automatisierter Erzeugung vereinfachter
    Steuermodelle via Machine Learning. Die dort gewählte
    Methodenfamilie (Regression / Random Forest) ist jedoch nicht
    physikalisch interpretierbar; \sindy{} und \edmd{} schließen diese
    Lücke.
\end{itemize}

Im Umfeld der Motorensystemidentifikation sind insbesondere die Arbeiten
der Yahagi-Gruppe einschlägig \cite{yahagiSparseIdentificationNonlinear,
yahagiSparseIdentificationNonlinear2025, yonezawaSparseIdentificationNonlinear2025}.
Diese Arbeiten setzen \sindy{} für die Identifikation einzelner
Untersysteme von PKW-Motoren ein und berichten kompakte, regelbasiert
einsetzbare Modelle. Eine systematische Übertragung auf
mittelschnelllaufende Großmotoren mit ihren um Größenordnungen längeren
thermischen Zeitkonstanten ist nicht beschrieben. Für die industrielle
Robustheit ist die Ensemble-Linie
\cite{faselEnsembleSINDyRobustSparse2022} sowie die auf industrielle
Prozesse zugeschnittene Sparse-Identifikation (EMEC-nahe Ansätze,
\cite{bhadrirajuOASISPOperableAdaptive2021}) einschlägig.

Für \edmd{} und Koopman-Methoden im Antriebsstrang sind die Beiträge von
\cite{medaKoopmanBasedMethodsEV2025} (Antriebsstrang von Elektrofahrzeugen)
und \cite{zinageDataDrivenModeling2022} exemplarisch. Beide bestätigen
das in Abschnitt~\ref{ssec.sot.edmd.def} skizzierte Bild, dass
\edmd-Modelle bei sorgfältig gewählter Liftung präzise Vorhersagen bei
moderaten Datenmengen erlauben. Für die modellprädiktive Regelung ist
\cite{kordaLinearPredictorsNonlinear2018} die Kernreferenz, weil die
globale lineare Repräsentation im gelifteten Raum die Formulierung des
Regelproblems als quadratisches Programm erlaubt
(Abschnitt~\ref{sec.methodik.mpc}); der
% ALT (Copilot 01.09.): \ref{sec.sot.mpc} war tot, das Label existiert nur
% noch in _alt\ch3_old.tex; die MPC-Formulierung steht jetzt in Kap. 5.
% 2. Durchgang nach Lektor: "klassische lineare MPC-Werkzeuge nutzbar" auf
% QP-Formulierung zurückgenommen, die eigene Umsetzung nutzt keine
% MPC-Toolbox (Lizenz), sondern einen MATLAB-Funktionsblock.
Stand der Koopman-MPC-Theorie ist in
\cite{strasserOverviewKoopmanbasedControl2026} konsolidiert, und
\cite{dahdahClosedloopKoopmanOperator2024} adressiert den
Closed-Loop-Bias auf Koopman-Seite.

Eine Anwendung von \sindy{} oder \edmd{} auf die Ladelufttemperatur eines
Dual-Fuel-Großmotors ist in der Sichtung für diese Arbeit nicht
aufgetreten. Grundlage sind der für die Arbeit aufgebaute Zotero-Bestand
und eine Websuche, beide mit Stand 02.09.2026, gesucht über die
Verfahrensnamen SINDy, SINDYc, Koopman, eDMD und eDMDc in Verbindung mit
Ladeluft, Ladeluftkühler, charge air, intercooler, Dual-Fuel und
Großmotor. Aus derselben Sichtung stammen zwei weitere Leerbefunde: keine
der gefundenen \sindy- oder Koopman-Arbeiten prüft ihre Residuen auf
Autokorrelation, Kreuzkorrelation mit dem Eingang oder Kohärenz, und ein
quasi-LPV-Prädiktor für die Ladelufttemperatur unter modellprädiktiver
Regelung ist nicht beschrieben. Alle drei Aussagen gelten für diesen
Rechercheumfang und sind kein Nachweis der Nichtexistenz.



%%%%%%%%%%%%%%%%%%%%%%%%%%%%%%%%%%%%%%%%%%%%%%%%%%%%%%%%%%%%%%%%%%%%%
%%  Einordnung und Abgrenzung der Arbeit
%%%%%%%%%%%%%%%%%%%%%%%%%%%%%%%%%%%%%%%%%%%%%%%%%%%%%%%%%%%%%%%%%%%%%



% COPILOT-ZUSATZ BEGIN ch3-15
\par\smallskip
\noindent\textbf{Copilot: Ergänzende Stichpunkte und Prüfhinweise}
\par\smallskip
\begin{itemize}
  \item \textbf{Großmotor-Kühlkreis.}
    Rupprecht: gekoppelte Ladeluft- und HT/LT-Kühlkreise mit
    physikalischen Komponentenmodellen; Temperaturregelung über
    die Ventilstellung im LT-Pfad.
    Parametrierung aus Theorie und Auslegungsdaten; Prüfung
    stationärer Punkte sowie offener und geschlossener Verläufe
    mit realen Prüfstandsdaten.
    Grenze: LT-Eintrittstemperatur als konstant angenommen;
    verbleibende Abweichungen am Ausgang des nachgebildeten
    vorhandenen Reglers trotz guter Wiedergabe der Ladelufttemperatur.
    Anschluss: Systemgrenze, Aktuator-/Sensorpfad und getrennte
    Prüfung von Temperatur und Stellsignal
    \cite[S.~12, 44--47, 96--100 und 138--140]{rupprechtAnalysisSimulationOptimisation2016}.
  \item \textbf{Reduzierte Kühlerdynamik.}
    Vagapov et al.: Hammerstein-Modell aus nichtlinearer
    stationärer Kühlerbeschreibung und linearer Dynamik,
    geprüft mit Prüfstands- und Fahrzeugmessungen.
    Serielle beziehungsweise parallele \emph{virtuelle} Wärmeübertrager
    berücksichtigen den Wärmeverlust an die Umgebung.
    Übertragbarer Baustein: Trennung von Statik und Dynamik,
    Ordnungsreduktion anhand der Beiträge zur Kühlerantwort.
    Grenze: Modellvorhersage, kein geschlossener Reglernachweis;
    eigene HT/LT-Anordnung, dominierende Speicher und Parameter
    gesondert begründen
    \cite[Abschn.~2--5]{vagapovDynamicModelTemperature2022}.
  \item \textbf{Prädiktive Ladelufttemperaturregelung im Fahrzeug.}
    Vagapov: äußere prädiktive Stufe für die Lufttemperatur nach dem
    Kühler; Vorgaben für Kühlmittelmassenstrom und -eintrittstemperatur,
    Umsetzung über Pumpe und Ventil.
    Die vollständige Kaskade erfüllt die geforderte Rechenzeit nicht;
    Fahrzeugversuche mit einer vereinfachten inneren modellbasierten
    Steuerung und verbleibender äußerer prädiktiver Stufe.
    Regelgüte dort nur grob untersucht, Gewichte und Nebenbedingungen
    beispielhaft gewählt.
    Anschluss: Modellgröße, Optimierungsaufwand und Reglerstruktur
    zusammen bewerten; Fahrversuch als Umsetzungsbeleg mit diesen
    Grenzen und den Anlagenunterschieden zum Großmotor
    \cite[Abschn.~9.1 und 9.4, insbesondere S.~137--140]{vagapovModellierungIdentifikationUnd2024}.
  \item \textbf{Kühler als Teil eines Gesamtmotormodells.}
    Sui et al.: stationäre Kühlerbeschreibung über Effektivität,
    Luftmassenstrom und Eintrittstemperaturen innerhalb eines
    Marine-Mittelwertmodells.
    Motorvalidierung anhand quasi-stationärer Prüfstandsdaten;
    kein isolierter dynamischer Nachweis des Kühlerbausteins.
    Anschluss: thermische Struktur der Eingangsmerkmale;
    dynamische Ventilverstärkung der eigenen Strecke zusätzlich prüfen
    \cite[Abschn.~4.1; Anhang~B.2, Gl.~(B.6)--(B.10)]{suiMeanValueFirst2022}.
\end{itemize}

\begin{itemize}
  \item \textbf{SINDy und NMPC.}
    Yahagi et al. modellieren Ladedruck und AGR-Rate mit einem
    diskreten spärlichen Dynamikmodell und untersuchen darauf
    aufbauend nichtlineare modellprädiktive Regelung.
    Stellgrößen: Turbinengeometrie und AGR-Ventil;
    Drehzahl und Einspritzmenge als weitere Eingänge.
    Identifikation und Reglervergleich am simulierten
    Mean-Value-Motormodell
    \cite[Abschn.~3 und 5]{yahagiSparseIdentificationNonlinear2025}.
  \item \textbf{Koopman-Prädiktion mit realen geregelten Messdaten.}
    Yahagi et al. verwenden Ein-/Ausgangshistorien und eine
    verallgemeinerte bilineare Koopman-Rekursion.
    Zusätzlich zur Simulation: Motorprüfstandsdaten unter
    Serien-PID-Regelung; betrachtete Ausgänge dort Ladedruck und
    Luftmassenstrom. Nachweis der Modellvorhersage;
    Entwicklung eines Koopman-Reglers als nächste Aufgabe.
    Quellenstand: arXiv-Preprint vom Februar 2026
    \cite[Abschn.~III, V und VI]{yahagiGeneralizedBilinearKoopman2026}.
  \item \textbf{Übertragbare methodische Frage.}
    In beiden Arbeiten: kompakte Dynamikmodelle für einen angeregten
    Motorluftpfad. Für diese Arbeit zusätzlich zu prüfen:
    Temperaturvorhersage und thermischer Ventilpfad.
    Die bilineare Koopman-Struktur zudem von der eigenen affinen
    Rekursion unterscheiden; eine günstigere Vorhersage in einer
    anderen Anlage ist keine allgemeine Rangfolge der Modellklassen.
    Bezug zu Abschnitt~\ref{ssec.sot.edmd.def}.
\end{itemize}
% COPILOT-ZUSATZ ENDE ch3-15
\section{Einordnung und Abgrenzung der Arbeit}
\label{sec.sot.einordnung}


Aus dem vorstehenden Überblick folgt: Eine Arbeit, die die
Ladelufttemperatur eines Großmotors für eine prädiktive Regelung
modelliert, findet sich in der gesichteten Literatur nicht. Die
einzelnen Bausteine sind dagegen vielfach behandelt.

Die dynamische Modellierung der Ladelufttemperatur von
Verbrennungsmotoren ist Teil vieler vorangehender Arbeiten, steht dort
aber selten im Fokus. Hervorzuheben ist
\cite{vagapovDynamicModelTemperature2022}, weil dort die
Ladelufttemperatur explizit behandelt wird und der Aktor wie im
vorliegenden Fall ein Dreiwegeventil ist.

Zur Anwendung von \sindyc{} und eDMDc existieren trotz der relativen
Neuheit von \sindy{} viele Arbeiten; nur wenige befassen sich mit
Verbrennungsmotoren \cite{yahagiSparseIdentificationNonlinear2025}. In
den in Abschnitt~\ref{sec.sot.application} gesichteten Arbeiten stammen
die Messdaten fast immer aus Simulationen. Eine Ausnahme mit
echten Messungen an einem Industriesystem ist
\cite{yahagiGeneralizedBilinearKoopman2026}. Für die hier verfolgte
Frage ist dieser Weg zirkulär, weil ein aus einem Modell gewonnenes
Modell die Güte seiner Quelle nicht übertreffen kann; die vorliegende
Arbeit arbeitet deshalb auf Messdaten (Kapitel~\ref{ch.datenbasis}).

Zur prädiktiven Regelung liegen viele Arbeiten vor, darunter
Anwendungen, in denen eine Temperatur geregelt wird, etwa
\cite{kelmanBilinearModelPredictive2011} für eine Gebäudeklimaanlage.
Die Ladelufttemperatur von Verbrennungsmotoren behandelt keine davon.


Vom rein physikalischen Modellieren grenzt sich die Arbeit klar ab.
Untersucht wird, wie sich mit \sindy{} und \edmd{} schnell und ohne
vollständige First-Principles-Parametrierung ein Regelmodell erstellen
lässt; die Physik geht als Auswahl der Bibliothek ein, nicht als fertig
parametriertes Modell.

Neuronale Netze (RNN, LSTM, NeuralODE, Transformer) sind eine
naheliegende dritte Familie für die nichtlineare Modellierung von
Zeitreihen. Sie werden in dieser Arbeit bewusst nicht als
Vergleichskandidaten aufgenommen, weil:

\begin{itemize}
  \item \textbf{Interpretierbarkeit.} Eine zentrale Anforderung dieser
    Arbeit ist die Lesbarkeit des identifizierten Modells in Termen
    physikalischer Größen, was Black-Box-NN nicht erfüllen.
  \item \textbf{Übertragbarkeit.} Die Übertragung trainierter NN zwischen
    Motorbaureihen würde aufwendiges Transfer-Lernen erfordern.
    Sparse-symbolische Modelle erlauben dagegen einen direkten
    Strukturvergleich.
\end{itemize}

% COPILOT-ZUSATZ BEGIN ch3-16
\par\smallskip
\noindent\textbf{Copilot: Ergänzende Stichpunkte und Prüfhinweise}
\par\smallskip
% COPILOT-STICHPUNKTE (09.09.2026; Literaturvertiefung; Kapitelort: Synthese)
% Eigene Ableitung aus den einzeln geprüften Arbeiten des vorigen Abschnitts.
% Modellrollen: 2026-09-08_Argumentation_Belege/MODELLE_VERSTEHEN.md:9-23,37-70,103-172.
% Vergleichsgrenzen: VERTRAG_279.md:3-12; METHODIK_VERGLEICH_20260908.md:28-48.
\begin{itemize}
  \item \textbf{Anwendungsbezogene Forschungsfrage:}
    Welche kompakte Modellstruktur beschreibt aus den verfügbaren
    Großmotor-Messdaten sowohl $T_{\mathrm{eng,in}}$ als auch die
    Wirkung des LT-Ventils über den vorgesehenen Regelhorizont?
    Physikalische Struktur, sparsame Parametrierung und geliftete
    Vorhersage als miteinander zu prüfende Möglichkeiten.
  \item \textbf{Anschluss an die Literatur:}
    Systemgrenze aus der Kühlkreismodellierung, stationäre
    Wärmeübertragung und reduzierte Speicherwirkung aus der
    thermischen Modellierung, rekursive Bewertung aus den
    datengetriebenen Luftpfadarbeiten.
    Der eigene Beitrag liegt in ihrer nachvollziehbaren Umsetzung
    und Prüfung für die konkrete Temperaturregelstrecke.
  \item \textbf{Getrennte Nachweise:}
    Parameterschätzung, Auswahl auf Entwicklungsdaten, unabhängige
    Modellbewertung und simulativer Reglervergleich.
    Retrospektive Vergleiche bereits bekannter Modellfassungen
    entsprechend kennzeichnen. Aufgezeichnete zukünftige
    Störgrößen bei Offline-Prognosen sind eine eigene
    Informationsannahme.
  \item \textbf{Reichweite:}
    explizite Terme und überschaubare Modellstrukturen als
    Entwurfsanforderung; Genauigkeit, Stellwirkung, Robustheit und
    Übertragbarkeit als Ergebnisse der jeweiligen Prüfung.
    Keine generelle Überlegenheit gegenüber neuronalen Verfahren
    und kein Nachweis der Nichtexistenz anderer passender Arbeiten
    aus dem begrenzten Literaturvergleich.
\end{itemize}
% COPILOT-ZUSATZ ENDE ch3-16
